{\sloppy\scriptsize\setlength{\tabcolsep}{3pt}
\begin{longtable}{>{\raggedright\arraybackslash}p{0.05\textwidth}>{\raggedright\arraybackslash}p{0.065\textwidth}>{\raggedright\arraybackslash}p{0.425\textwidth}>{\raggedright\arraybackslash}p{0.24\textwidth}>{\raggedright\arraybackslash}p{0.1\textwidth}}
\caption{The 52 predictions whose runs the main text or this supplement reports, of the 98 the study registered: 34 held and 18 failed, 6 of the failures foreseen. Each is given as its predictions file states it (shortened where long), with the expectation recorded before the runs and the verdict. Labels are defined in Section~\ref{sec:s-key}.}\label{tab:s-registry}\\
\toprule
Step & & Prediction & Expectation & Verdict \\
\midrule
\endfirsthead
\multicolumn{5}{l}{\textit{Continued}} \\
\toprule
Step & & Prediction & Expectation & Verdict \\
\midrule
\endhead
5 & P1 & At the central contact conductance with the posterior-mean kinetics, the resolved neighbour propagates in series and parallel, its time to propagation within the measured 147 to 342 s. & Study's expectation: May fail: by the estimates above the hot spot stays below 255.7\degC, or, if the line of contact rises further, reaches it within the first minute. & failed \\
5 & P2 & In both connections, series and parallel, the lumped neighbour propagates more than 20 s later than the resolved one, or not at all within 900 s. & Study's expectation: Held, most likely as "not at all": its passive peak should lie near 145 to 170\degC. & held \\
5 & P3 & At the resolved model's propagation time, its mid-surface temperature lies between 150 and 255.7\,\textdegree C: the hot spot leads the surface a thermocouple sees. & Study's expectation: Toward the trigger, held if propagation occurs (that face is the line of contact); away from it, may fail below 150\degC. & failed \\
5 & P4 & The resolved model's mean mid-surface heating rate over the 60 s before propagation lies between 0.2 and 0.8\,\textdegree C/s (measured about 0.4). & Study's expectation: Uncertain. & failed \\
5 & P5 & The predicted series and parallel times to propagation differ by less than 30 s (the measured means differ by 8 s). & Study's expectation: Held if both propagate: the tabs carry 4 to 8 \% of the contact path. & failed \\
5 & P6 & The share of the 4,800 factorial runs in which the resolved neighbour propagates lies between 30 and 90 \% (measured 9 of 16, 56 \%). & Study's expectation: May fail low: the model lacks the vent, ejecta and flame heat, and the estimates above fall short of 255.7\degC{} at the central contact. & failed \\
5 & P7 & The contact and tab conductance range carries more of the spread in predicted time to propagation than the kinetic posterior (standard deviation ratio above 1), reversing step 3's onset finding. & Study's expectation: Held if enough runs propagate to give both spreads: the contact range changes the path 4.6-fold. & failed \\
7 & P1 & At the central contact with per-test triggers, C1 (hot spot at 255.7\degC) predicts propagation in fewer than 9 of the 16 tests. & Plan's expectation: Likely. Study's expectation: Holds, with eight. & held \\
8 & P1 & Of the nine propagated tests, $f^*$(C1) is 0 in the six step 7 hits and below 0.3 in the other three. & Plan's expectation: Likely. Study's expectation: Holds. & held \\
8 & P2 & Of the seven non-propagated tests, at least three, including the two step 7 false alarms at 0, have $f^*$(C1) below 0.3, so the propagated and non-propagated tests' $f^*$(C1) ranges overlap. & Plan's expectation: Likely. Study's expectation: Holds, with all seven below 0.3. & held \\
8 & P3 & $f^*$(C2) lies between 0.05 and 0.5 in at least twelve of the sixteen tests. & Plan's expectation: Uncertain. Study's expectation: Fails. & failed, foreseen \\
8 & P6 & Within each connection, $f^*$(C1) and $f^*$(C2) never increase with trigger energy (no inversion). & Plan's expectation: Likely; a check on the mapping. Study's expectation: Holds, with ties. & held \\
9 & P1 & With the 21.39 s delivery, C1 reproduces no propagated test's measured time within 30 s at any fraction at or above $f^*$'s lower end. & Plan's expectation: Likely. Study's expectation: Holds. & held \\
9 & P2 & With the 21.39 s delivery, C2 reproduces at least six of nine propagation times within 30 s, at a fraction between $f^*$'s lower end and 0.1 above its upper end. & Plan's expectation: Uncertain. Study's expectation: Holds. & held \\
9 & P3 & Every C2 time-matching fraction found in P2 lies no more than 0.02 above the upper end of $f^*$(C2). & Plan's expectation: Uncertain. Study's expectation: Uncertain; it turns on PP-Exp1. & held \\
9 & P4 & For the five low-maximum non-propagated tests, the peak-matching fraction at the thermocouple node lies below both $f^*$(C1) and $f^*$(C2). & Plan's expectation: Likely. Study's expectation: Fails, on ST-Exp1 and ST-Exp2, whatever the runs show. & failed, foreseen \\
9 & P5 & The peak-matching fractions (five non-propagated tests) and the C2 time-matching fractions (nine propagated tests) separate: the largest of the former is below the smallest of the latter. & Plan's expectation: Uncertain; this is the step's question. Study's expectation: Fails. & failed, foreseen \\
9 & P7 & The bracket on the ejected-to-body energy ratio rescales every time- and peak-matching fraction by 0.75 to 1.40 and changes P5's verdict in neither direction. & Plan's expectation: Likely. Study's expectation: Holds. & held \\
10 & P1 & Before propagation, the model under the 21.39 s delivery, at one end of its threshold bracket, fits PT-Exp4's measured history at the step 9 (far-face) node within 25 K root-mean-square. & Plan's expectation: uncertain. & held \\
10 & P6 & Held out, no predictor beats one probability for every test by more than 0.1 in mean log score in either split direction, and no permutation p-value is below 0.05. & Plan's expectation: likely. & held \\
10 & P8 & The survey finds at least one dataset with per-test outcomes and times for cells in direct contact, and none whose cell this study's parameters cover without an invented value. & Plan's expectation: uncertain. & failed \\
11 & P1 & For each of the four tests, $f^*$(C2) under every kinetic draw lies within a factor of two of its posterior-mean value, or within 0.025 of it where that is 0. & Plan's expectation: likely. & failed \\
11 & P2 & For each of the four tests, $f^*$(C1) under every kinetic draw lies within a factor of two of its posterior-mean value, or within 0.025 of it where that is 0. & Plan's expectation: uncertain; the hot spot's onset is local and the kinetics set it. & held \\
11 & P3 & Under each of the eight draws, the range of the C2 threshold energy across the four tests exceeds the largest within-test band of that energy across the draws. & Plan's expectation: likely. & held \\
11 & P4 & In the four tests with zero peak-matching fraction, the far face with no delivery exceeds the measured maximum under at least three of four thermal draws and at both contacts. & Plan's expectation: uncertain; the low contact and the smallest radial conductivity are the two that could remove it; the predictions file's analysis: “P4, which asks that it do so at both contacts, fails before any step 11 run” & failed, foreseen \\
11 & P5 & In each of the four tests the band of $f^*$(C2) across the four thermal draws is narrower than its band across the four kinetic draws. & Plan's expectation: uncertain. & failed \\
12 & P1 & With $E_\mathrm{b}$ capped at 33.2 kJ, conduction and radiation alone meet C1 in none of the eight tests at the central contact. & Plan's expectation: likely (step 5 found no ignition at 33.2 kJ with the two-cell study's trigger); the predictions file records it as fixed by step 6's records (Section~\ref{sec:s-registry}). & held \\
12 & P2 & At the high contact the same holds. & Plan's expectation: uncertain; the predictions file records it as close to fixed (step 6, with other tabs; Section~\ref{sec:s-registry}). & held \\
12 & P3 & With $E_\mathrm{e}$ fixed at the Databank's median, the delivered energies $E^*$ across the nine propagated tests still span at least tenfold. & Plan's expectation: likely; the predictions file records it as fixed by step 9's records. & held \\
12 & P4 & With $E_\mathrm{e}$ fixed, the propagated and non-propagated tests' fractions still overlap (the largest $f_\mathrm{max}$ above the smallest $f_\mathrm{t}$(C2)). & Plan's expectation: likely; the predictions file records it as fixed by step 9's records. & held \\
13 & P1 & For PT-Exp1 with a 5 s triangular delivery over the facing half, the resolved $f^*$(C1) lies within 10 \% of the skin model's closed form (0.07701). & likely & held \\
13 & P2 & At 40 s (PT-Exp1, facing half), the cross-section conduction model (0.21775) lies within 5 \% of the resolved bracket's midpoint. & uncertain & held \\
13 & P3 & At 120 s (PT-Exp1, facing half), the resolved threshold lies more than 5 \% below the cross-section conduction model's (0.41450): local ignition. & uncertain & failed \\
13 & P4 & For ST-Exp3 and PT-Exp3, the cross-section model with the trigger's conduction lies within 10 \% of the resolved midpoints at 10 and 60 s, and gives the direction of change from 21.39 s in all four cases. & uncertain & held \\
13 & P7 & The ignition temperature of the posterior-mean kinetics lies between 215 and 240\degC{} (known when registered, so not scored), and the central 95 \% of the 200 draws spans less than 40 K (scored). & likely for the mean, uncertain for the span & held \\
13 & P8 & The two-body model's whole-cell threshold lies within the larger of 10 \% and 300 J of the lumped model's midpoint in at least 28 of the 32 low and high contact searches. & uncertain & held \\
13 & P9 & For PT-Exp1, the ratio $f^*$(C2)/$f^*$(C1) falls from 21.39 to 120 to 360 s and is below 1.5 at 360 s. & likely for the fall, uncertain for the level & held \\
13 & P10 & For PP-Exp3, ST-Exp2 and PT-Exp3, the lumped model's needed delivery falls by 150 to 300 J per kJ of the trigger's body energy across the 4 kJ between shifted runs. & likely for the first two, uncertain for PT-Exp3 & failed \\
14 & P1 & The upper bound of radiation from a cap-end vent to the facing half is below 2 \% of 76.74 kJ. & the plan's line for it says likely; the predictions file's analysis, written before any bound, says P1 is expected to fail & failed, foreseen \\
14 & P2 & The upper bounds of the cap-end paths that are determinable sum to less than 5.96 kJ, the largest need among the seven matched times. & uncertain; the predictions file's analysis says P2 is expected to fail & failed, foreseen \\
14 & P3 & They sum to more than 0.26 kJ, the smallest need. & likely & held \\
14 & P4 & The upper bound of a breach's jet facing the neighbour exceeds 5.96 kJ delivered within 21.39 s. & likely & held \\
14 & P5 & At least one of the six paths is not determinable from published inputs for this cell. & likely & held \\
14 & P6 & For the parallel tests the circuit's upper bound on the neighbour's own heating exceeds 1 kJ. & uncertain; the predictions file's analysis says P6 is expected to hold & held \\
15 & P1 & The reactive cross-section's $E^*$(C8) lies within 15 \% of the midpoint of the three-dimensional bracket in at least three of four tests. & uncertain; the reduced model's values are the searches' centres & held \\
15 & P2 & For PT-Exp4 with the 21.39 s delivery, $E^*$(C8) lies between 6 and 10 kJ. & likely; the reduced model's value, registered with it: 7,062 J & held \\
15 & P3 & In at least three of the four tests $E^*$(C8) is at least 1.25 times the whole-cell clock's threshold $E^*$(C2). & likely for the three strong triggers, uncertain for ST-Exp3; the reduced model's values: 1.49, 1.72, above zero in SP-Exp3, 0.85 & held \\
15 & P4 & For PT-Exp4, in the run at the upper end of the bracket, the mean crosses 255.7\degC{} between 150 and 400 s after the trigger's start. & likely; the reduced model's value: 310.6 s at the upper end of its own bracket & held \\
15 & P5 & In that run the far face reads more than 15 K above 197.3\degC, the reading from which the measured neighbour jumped, when its own jump begins. & likely; the reduced model's value: 240.1\,\textdegree C & held \\
15 & P6 & The constant power after the 5.0 kJ pulse that brings the cell's own runaway to 271.5 s lies between 5 and 25 W. & likely; the reduced model's value: 12.5 W & held \\
15 & P7 & Under none of four kinetic draws and not with the aged SEI does the resolved cell run away at PT-Exp4's 5,045 J, and the least shift of the rate laws that makes it run away is 10 K or more. & likely; the reduced model's values: no runaway in any of the five runs, the least shift 20.0 K & held \\
15 & P8 & PT-Exp4's run with 6,209 J over 240 s, carried to 900 s, runs away. & likely; the reduced model's value, registered with it: no runaway, the mean peaking at 220.5\,\textdegree C & failed \\
\bottomrule
\end{longtable}
}
